% Options for packages loaded elsewhere
\PassOptionsToPackage{unicode}{hyperref}
\PassOptionsToPackage{hyphens}{url}
%
\documentclass[
]{article}
\usepackage{amsmath,amssymb}
\usepackage{iftex}
\ifPDFTeX
  \usepackage[T1]{fontenc}
  \usepackage[utf8]{inputenc}
  \usepackage{textcomp} % provide euro and other symbols
\else % if luatex or xetex
  \usepackage{unicode-math} % this also loads fontspec
  \defaultfontfeatures{Scale=MatchLowercase}
  \defaultfontfeatures[\rmfamily]{Ligatures=TeX,Scale=1}
\fi
\usepackage{lmodern}
\ifPDFTeX\else
  % xetex/luatex font selection
\fi
% Use upquote if available, for straight quotes in verbatim environments
\IfFileExists{upquote.sty}{\usepackage{upquote}}{}
\IfFileExists{microtype.sty}{% use microtype if available
  \usepackage[]{microtype}
  \UseMicrotypeSet[protrusion]{basicmath} % disable protrusion for tt fonts
}{}
\makeatletter
\@ifundefined{KOMAClassName}{% if non-KOMA class
  \IfFileExists{parskip.sty}{%
    \usepackage{parskip}
  }{% else
    \setlength{\parindent}{0pt}
    \setlength{\parskip}{6pt plus 2pt minus 1pt}}
}{% if KOMA class
  \KOMAoptions{parskip=half}}
\makeatother
\usepackage{xcolor}
\usepackage[margin=1in]{geometry}
\usepackage{color}
\usepackage{fancyvrb}
\newcommand{\VerbBar}{|}
\newcommand{\VERB}{\Verb[commandchars=\\\{\}]}
\DefineVerbatimEnvironment{Highlighting}{Verbatim}{commandchars=\\\{\}}
% Add ',fontsize=\small' for more characters per line
\usepackage{framed}
\definecolor{shadecolor}{RGB}{248,248,248}
\newenvironment{Shaded}{\begin{snugshade}}{\end{snugshade}}
\newcommand{\AlertTok}[1]{\textcolor[rgb]{0.94,0.16,0.16}{#1}}
\newcommand{\AnnotationTok}[1]{\textcolor[rgb]{0.56,0.35,0.01}{\textbf{\textit{#1}}}}
\newcommand{\AttributeTok}[1]{\textcolor[rgb]{0.13,0.29,0.53}{#1}}
\newcommand{\BaseNTok}[1]{\textcolor[rgb]{0.00,0.00,0.81}{#1}}
\newcommand{\BuiltInTok}[1]{#1}
\newcommand{\CharTok}[1]{\textcolor[rgb]{0.31,0.60,0.02}{#1}}
\newcommand{\CommentTok}[1]{\textcolor[rgb]{0.56,0.35,0.01}{\textit{#1}}}
\newcommand{\CommentVarTok}[1]{\textcolor[rgb]{0.56,0.35,0.01}{\textbf{\textit{#1}}}}
\newcommand{\ConstantTok}[1]{\textcolor[rgb]{0.56,0.35,0.01}{#1}}
\newcommand{\ControlFlowTok}[1]{\textcolor[rgb]{0.13,0.29,0.53}{\textbf{#1}}}
\newcommand{\DataTypeTok}[1]{\textcolor[rgb]{0.13,0.29,0.53}{#1}}
\newcommand{\DecValTok}[1]{\textcolor[rgb]{0.00,0.00,0.81}{#1}}
\newcommand{\DocumentationTok}[1]{\textcolor[rgb]{0.56,0.35,0.01}{\textbf{\textit{#1}}}}
\newcommand{\ErrorTok}[1]{\textcolor[rgb]{0.64,0.00,0.00}{\textbf{#1}}}
\newcommand{\ExtensionTok}[1]{#1}
\newcommand{\FloatTok}[1]{\textcolor[rgb]{0.00,0.00,0.81}{#1}}
\newcommand{\FunctionTok}[1]{\textcolor[rgb]{0.13,0.29,0.53}{\textbf{#1}}}
\newcommand{\ImportTok}[1]{#1}
\newcommand{\InformationTok}[1]{\textcolor[rgb]{0.56,0.35,0.01}{\textbf{\textit{#1}}}}
\newcommand{\KeywordTok}[1]{\textcolor[rgb]{0.13,0.29,0.53}{\textbf{#1}}}
\newcommand{\NormalTok}[1]{#1}
\newcommand{\OperatorTok}[1]{\textcolor[rgb]{0.81,0.36,0.00}{\textbf{#1}}}
\newcommand{\OtherTok}[1]{\textcolor[rgb]{0.56,0.35,0.01}{#1}}
\newcommand{\PreprocessorTok}[1]{\textcolor[rgb]{0.56,0.35,0.01}{\textit{#1}}}
\newcommand{\RegionMarkerTok}[1]{#1}
\newcommand{\SpecialCharTok}[1]{\textcolor[rgb]{0.81,0.36,0.00}{\textbf{#1}}}
\newcommand{\SpecialStringTok}[1]{\textcolor[rgb]{0.31,0.60,0.02}{#1}}
\newcommand{\StringTok}[1]{\textcolor[rgb]{0.31,0.60,0.02}{#1}}
\newcommand{\VariableTok}[1]{\textcolor[rgb]{0.00,0.00,0.00}{#1}}
\newcommand{\VerbatimStringTok}[1]{\textcolor[rgb]{0.31,0.60,0.02}{#1}}
\newcommand{\WarningTok}[1]{\textcolor[rgb]{0.56,0.35,0.01}{\textbf{\textit{#1}}}}
\usepackage{graphicx}
\makeatletter
\def\maxwidth{\ifdim\Gin@nat@width>\linewidth\linewidth\else\Gin@nat@width\fi}
\def\maxheight{\ifdim\Gin@nat@height>\textheight\textheight\else\Gin@nat@height\fi}
\makeatother
% Scale images if necessary, so that they will not overflow the page
% margins by default, and it is still possible to overwrite the defaults
% using explicit options in \includegraphics[width, height, ...]{}
\setkeys{Gin}{width=\maxwidth,height=\maxheight,keepaspectratio}
% Set default figure placement to htbp
\makeatletter
\def\fps@figure{htbp}
\makeatother
\setlength{\emergencystretch}{3em} % prevent overfull lines
\providecommand{\tightlist}{%
  \setlength{\itemsep}{0pt}\setlength{\parskip}{0pt}}
\setcounter{secnumdepth}{-\maxdimen} % remove section numbering
\ifLuaTeX
  \usepackage{selnolig}  % disable illegal ligatures
\fi
\usepackage{bookmark}
\IfFileExists{xurl.sty}{\usepackage{xurl}}{} % add URL line breaks if available
\urlstyle{same}
\hypersetup{
  pdftitle={Danish Kings, W10 assignment},
  pdfauthor={Christian Reinholdt},
  hidelinks,
  pdfcreator={LaTeX via pandoc}}

\title{Danish Kings, W10 assignment}
\author{Christian Reinholdt}
\date{16.5.2025}

\begin{document}
\maketitle

\begin{Shaded}
\begin{Highlighting}[]
\NormalTok{knitr}\SpecialCharTok{::}\NormalTok{opts\_chunk}\SpecialCharTok{$}\FunctionTok{set}\NormalTok{(}\AttributeTok{echo =} \ConstantTok{TRUE}\NormalTok{)}
\end{Highlighting}
\end{Shaded}

\subsection{Installing some packages first (change eval=TRUE, if
packages are
needed)}\label{installing-some-packages-first-change-evaltrue-if-packages-are-needed}

\begin{Shaded}
\begin{Highlighting}[]
\FunctionTok{options}\NormalTok{(}\AttributeTok{repos =} \FunctionTok{c}\NormalTok{(}\AttributeTok{CRAN =} \StringTok{"https://cloud.r{-}project.org"}\NormalTok{))}
\FunctionTok{install.packages}\NormalTok{(}\StringTok{"tidyverse"}\NormalTok{)}
\FunctionTok{install.packages}\NormalTok{(}\StringTok{"tidyverse"}\NormalTok{)}
\FunctionTok{install.packages}\NormalTok{(}\StringTok{"dplyr"}\NormalTok{)}
\FunctionTok{install.packages}\NormalTok{(}\StringTok{"readr"}\NormalTok{)}
\FunctionTok{install.packages}\NormalTok{(}\StringTok{"ggplot2"}\NormalTok{)}
\end{Highlighting}
\end{Shaded}

\subsection{When packages are already
installed:}\label{when-packages-are-already-installed}

\begin{Shaded}
\begin{Highlighting}[]
\FunctionTok{library}\NormalTok{ (}\StringTok{"tidyverse"}\NormalTok{)}
\end{Highlighting}
\end{Shaded}

\begin{verbatim}
## -- Attaching core tidyverse packages ------------------------ tidyverse 2.0.0 --
## v dplyr     1.1.4     v readr     2.1.5
## v forcats   1.0.0     v stringr   1.5.1
## v ggplot2   3.5.2     v tibble    3.2.1
## v lubridate 1.9.4     v tidyr     1.3.1
## v purrr     1.0.4     
## -- Conflicts ------------------------------------------ tidyverse_conflicts() --
## x dplyr::filter() masks stats::filter()
## x dplyr::lag()    masks stats::lag()
## i Use the conflicted package (<http://conflicted.r-lib.org/>) to force all conflicts to become errors
\end{verbatim}

\begin{Shaded}
\begin{Highlighting}[]
\FunctionTok{library}\NormalTok{ (}\StringTok{"dplyr"}\NormalTok{)}
\FunctionTok{library}\NormalTok{ (}\StringTok{"readr"}\NormalTok{)}
\FunctionTok{library}\NormalTok{ (}\StringTok{"ggplot2"}\NormalTok{)}
\end{Highlighting}
\end{Shaded}

The task here is to load your Danish Monarchs csv into R using the
\texttt{tidyverse} toolkit, calculate and explore the kings' duration of
reign with pipes \texttt{\%\textgreater{}\%} in \texttt{dplyr} and plot
it over time.

\subsection{Load the kings}\label{load-the-kings}

Make sure to first create an \texttt{.Rproj} workspace with a
\texttt{data/} folder where you place either your own dataset or the
provided \texttt{kings.csv} dataset.

\begin{enumerate}
\def\labelenumi{\arabic{enumi}.}
\tightlist
\item
  Look at the dataset that are you loading and check what its columns
  are separated by? (hint: open it in plain text editor to see)
\end{enumerate}

List what is the separator: commas (since it is a comma separated values
file (csv-file))

\begin{enumerate}
\def\labelenumi{\arabic{enumi}.}
\setcounter{enumi}{1}
\tightlist
\item
  Create a \texttt{kings} object in R with the different functions below
  and inspect the different outputs.
\end{enumerate}

\begin{itemize}
\tightlist
\item
  \texttt{read.csv()}
\item
  \texttt{read\_csv()}
\item
  \texttt{read.csv2()}
\item
  \texttt{read\_csv2()}
\end{itemize}

\begin{Shaded}
\begin{Highlighting}[]
\CommentTok{\# FILL IN THE CODE BELOW and review the outputs}

\NormalTok{kings1 }\OtherTok{\textless{}{-}} \FunctionTok{read.csv}\NormalTok{(}\StringTok{"data/kings data.csv"}\NormalTok{)}

\NormalTok{kings2 }\OtherTok{\textless{}{-}} \FunctionTok{read\_csv}\NormalTok{(}\StringTok{"data/kings data.csv"}\NormalTok{)}
\end{Highlighting}
\end{Shaded}

\begin{verbatim}
## Rows: 57 Columns: 11
## -- Column specification --------------------------------------------------------
## Delimiter: ","
## chr (11): Name, House, Start_year, End_year, Birth_year, Death_year, Gender,...
## 
## i Use `spec()` to retrieve the full column specification for this data.
## i Specify the column types or set `show_col_types = FALSE` to quiet this message.
\end{verbatim}

\begin{Shaded}
\begin{Highlighting}[]
\NormalTok{kings3 }\OtherTok{\textless{}{-}} \FunctionTok{read.csv2}\NormalTok{(}\StringTok{"data/kings data.csv"}\NormalTok{)}

\NormalTok{kings4 }\OtherTok{\textless{}{-}} \FunctionTok{read\_csv2}\NormalTok{( }\StringTok{"data/kings data.csv"}\NormalTok{)}
\end{Highlighting}
\end{Shaded}

\begin{verbatim}
## i Using "','" as decimal and "'.'" as grouping mark. Use `read_delim()` for more control.
\end{verbatim}

\begin{verbatim}
## Warning: One or more parsing issues, call `problems()` on your data frame for details,
## e.g.:
##   dat <- vroom(...)
##   problems(dat)
\end{verbatim}

\begin{verbatim}
## Rows: 57 Columns: 1
## -- Column specification --------------------------------------------------------
## Delimiter: ";"
## chr (1): Name,House,Start_year,End_year,Birth_year,Death_year,Gender,Dynasty...
## 
## i Use `spec()` to retrieve the full column specification for this data.
## i Specify the column types or set `show_col_types = FALSE` to quiet this message.
\end{verbatim}

Answer: 1. Which of these functions is a \texttt{tidyverse} function?
Read data with it below into a \texttt{kings} object - ``\_'' and ``.''
is a tidyverse function.

\begin{Shaded}
\begin{Highlighting}[]
\NormalTok{kings }\OtherTok{\textless{}{-}} \FunctionTok{read\_csv}\NormalTok{(}\StringTok{"data/kings data.csv"}\NormalTok{)}
\end{Highlighting}
\end{Shaded}

\begin{verbatim}
## Rows: 57 Columns: 11
## -- Column specification --------------------------------------------------------
## Delimiter: ","
## chr (11): Name, House, Start_year, End_year, Birth_year, Death_year, Gender,...
## 
## i Use `spec()` to retrieve the full column specification for this data.
## i Specify the column types or set `show_col_types = FALSE` to quiet this message.
\end{verbatim}

\begin{enumerate}
\def\labelenumi{\arabic{enumi}.}
\setcounter{enumi}{1}
\tightlist
\item
  What is the result of running \texttt{class()} on the \texttt{kings}
  object created with a tidyverse function.
\end{enumerate}

\begin{Shaded}
\begin{Highlighting}[]
\FunctionTok{class}\NormalTok{(kings)}
\end{Highlighting}
\end{Shaded}

\begin{verbatim}
## [1] "spec_tbl_df" "tbl_df"      "tbl"         "data.frame"
\end{verbatim}

\begin{itemize}
\tightlist
\item
  We see that ``class()'' informs us that the ``kings'' object is a
  tabular data.frame.
\end{itemize}

\begin{enumerate}
\def\labelenumi{\arabic{enumi}.}
\setcounter{enumi}{2}
\tightlist
\item
  How many columns does the object have when created with these
  different functions?
\end{enumerate}

\begin{itemize}
\tightlist
\item
  It differs, when using a proper tidyverse function to load the csv,
  the object has 11 columns. But when loading the csv improperly, it
  only has 1 column.
\end{itemize}

\begin{enumerate}
\def\labelenumi{\arabic{enumi}.}
\setcounter{enumi}{3}
\tightlist
\item
  Show the dataset so that we can see how R interprets each column
\end{enumerate}

\begin{Shaded}
\begin{Highlighting}[]
\FunctionTok{glimpse}\NormalTok{(kings)}
\end{Highlighting}
\end{Shaded}

\begin{verbatim}
## Rows: 57
## Columns: 11
## $ Name       <chr> "Gorm den Gamle", "Harald 1. Blåtand", "Toke_Gormsen", "Sve~
## $ House      <chr> "Gorm", "Gorm", "Gorm", "Gorm", "Gorm", "Gorm", "Gorm", "Fa~
## $ Start_year <chr> "936", "NULL", "985", "NULL", "1014", "1018", "1035", "1042~
## $ End_year   <chr> "958", "NULL", "986", "NULL", "1018", "1035", "1042", "1047~
## $ Birth_year <chr> "908", "936", "NULL", "NULL", "963", "NULL", "995", "1018",~
## $ Death_year <chr> "958", "985", "986", "1014", "NULL", "1035", "1042", "1047"~
## $ Gender     <chr> "M", "M", "M", "M", "M", "M", "M", "M", "M", "M", "M", "M",~
## $ Dynasty    <chr> "Jellingdynastiet", "Jellingdynastiet", "Jellingdynastiet",~
## $ Source     <chr> "https://da.wikipedia.org/wiki/Gorm_den_Gamle", "https://da~
## $ BirthDMY   <chr> "908", "936", NA, "17/04/963", "994", "995", "1018", "1024"~
## $ DeathDMY   <chr> "958", "987", NA, "03/02/1014", "1018", "12/11/1035", "08/0~
\end{verbatim}

We see that R interprets each column as ``characters'', that is, as
letters, and not for instance as numbers or years.

\subsection{Calculate the duration of reign for all the kings in your
table}\label{calculate-the-duration-of-reign-for-all-the-kings-in-your-table}

You can calculate the duration of reign in years with \texttt{mutate}
function by subtracting the equivalents of your \texttt{startReign} from
\texttt{endReign} columns and writing the result to a new column called
\texttt{duration}. But first you need to check a few things:

\begin{itemize}
\tightlist
\item
  Is your data messy? Fix it before re-importing to R
\item
  Data is not messy, since it is Adelas :-) I will however change how
  Rstudio interprets these columns, to be numeric.
\end{itemize}

\begin{Shaded}
\begin{Highlighting}[]
\NormalTok{kings }\OtherTok{\textless{}{-}}\NormalTok{ kings }\SpecialCharTok{\%\textgreater{}\%}
  \FunctionTok{mutate}\NormalTok{(}
    \AttributeTok{Name =} \FunctionTok{as.character}\NormalTok{(Name),}
    \AttributeTok{House =} \FunctionTok{as.character}\NormalTok{(House),}
    \AttributeTok{Start\_year =} \FunctionTok{as.numeric}\NormalTok{(Start\_year),}
    \AttributeTok{End\_year =} \FunctionTok{as.numeric}\NormalTok{(End\_year),}
    \AttributeTok{Birth\_year =} \FunctionTok{as.numeric}\NormalTok{(Birth\_year),}
    \AttributeTok{Death\_year =} \FunctionTok{as.numeric}\NormalTok{(Death\_year),}
    \AttributeTok{Gender =} \FunctionTok{as.character}\NormalTok{(Gender),       }
    \AttributeTok{Dynasty =} \FunctionTok{as.character}\NormalTok{(Dynasty),}
    \AttributeTok{Source =} \FunctionTok{as.character}\NormalTok{(Source),}
    \AttributeTok{BirthDMY =} \FunctionTok{as.character}\NormalTok{(BirthDMY),   }
    \AttributeTok{DeathDMY =} \FunctionTok{as.character}\NormalTok{(DeathDMY)}
\NormalTok{  ) }\SpecialCharTok{\%\textgreater{}\%}
  \FunctionTok{filter}\NormalTok{(}
    \SpecialCharTok{!}\FunctionTok{is.na}\NormalTok{(Start\_year),}
    \SpecialCharTok{!}\FunctionTok{is.na}\NormalTok{(End\_year)}
\NormalTok{  )}
\end{Highlighting}
\end{Shaded}

\begin{verbatim}
## Warning: There were 4 warnings in `mutate()`.
## The first warning was:
## i In argument: `Start_year = as.numeric(Start_year)`.
## Caused by warning:
## ! NAs introduced by coercion
## i Run `dplyr::last_dplyr_warnings()` to see the 3 remaining warnings.
\end{verbatim}

\begin{itemize}
\tightlist
\item
  Do your start and end of reign columns contain NAs? Choose the right
  strategy to deal with them: \texttt{na.omit()}, \texttt{na.rm=TRUE},
  \texttt{!is.na()} Create a new column called \texttt{duration} in the
  kings dataset, utilizing the \texttt{mutate()} function from
  tidyverse. Check with your group to brainstorm the options.
\end{itemize}

\begin{Shaded}
\begin{Highlighting}[]
\NormalTok{kings\_cleaned }\OtherTok{\textless{}{-}}\NormalTok{ kings }\SpecialCharTok{\%\textgreater{}\%}
  \FunctionTok{mutate}\NormalTok{(}\AttributeTok{duration =}\NormalTok{ End\_year }\SpecialCharTok{{-}}\NormalTok{ Start\_year) }
\end{Highlighting}
\end{Shaded}

\subsection{Calculate the average duration of reign for all
rulers}\label{calculate-the-average-duration-of-reign-for-all-rulers}

Do you remember how to calculate an average on a vector object? If not,
review the last two lessons and remember that a column is basically a
vector. So you need to subset your \texttt{kings} dataset to the
\texttt{duration} column. If you subset it as a vector you can calculate
average on it with \texttt{mean()} base-R function. If you subset it as
a tibble, you can calculate average on it with \texttt{summarize()}
tidyverse function. Try both ways!

\begin{itemize}
\tightlist
\item
  You first need to know how to select the relevant \texttt{duration}
  column. What are your options?
\item
  Is your selected \texttt{duration} column a tibble or a vector? The
  \texttt{mean()} function can only be run on a vector. The
  \texttt{summarize()} function works on a tibble.
\item
  Are you getting an error that there are characters in your column?
  Coerce your data to numbers with \texttt{as.numeric()}.
\item
  Remember to handle NAs: \texttt{mean(X,\ na.rm=TRUE)}
\end{itemize}

\begin{Shaded}
\begin{Highlighting}[]
\FunctionTok{library}\NormalTok{(dplyr)}

\NormalTok{mean\_duration }\OtherTok{\textless{}{-}}\NormalTok{ kings\_cleaned }\SpecialCharTok{\%\textgreater{}\%}
  \FunctionTok{summarize}\NormalTok{(}\AttributeTok{avg\_duration =} \FunctionTok{mean}\NormalTok{(duration, }\AttributeTok{na.rm =} \ConstantTok{TRUE}\NormalTok{))}

\NormalTok{mean\_duration}
\end{Highlighting}
\end{Shaded}

\begin{verbatim}
## # A tibble: 1 x 1
##   avg_duration
##          <dbl>
## 1         19.6
\end{verbatim}

\subsection{How many and which kings enjoyed a longer-than-average
duration of
reign?}\label{how-many-and-which-kings-enjoyed-a-longer-than-average-duration-of-reign}

You have calculated the average duration above. Use it now to
\texttt{filter()} the \texttt{duration} column in \texttt{kings}
dataset. Display the result and also count the resulting rows with
\texttt{count()}

\begin{Shaded}
\begin{Highlighting}[]
\NormalTok{long\_reign\_kings }\OtherTok{\textless{}{-}}\NormalTok{ kings\_cleaned }\SpecialCharTok{\%\textgreater{}\%}
  \FunctionTok{filter}\NormalTok{(duration }\SpecialCharTok{\textgreater{}}\NormalTok{ mean\_duration)}
\end{Highlighting}
\end{Shaded}

\begin{verbatim}
## Warning: Using one column matrices in `filter()` was deprecated in dplyr 1.1.0.
## i Please use one dimensional logical vectors instead.
## This warning is displayed once every 8 hours.
## Call `lifecycle::last_lifecycle_warnings()` to see where this warning was
## generated.
\end{verbatim}

\begin{Shaded}
\begin{Highlighting}[]
\FunctionTok{print}\NormalTok{(long\_reign\_kings)}
\end{Highlighting}
\end{Shaded}

\begin{verbatim}
## # A tibble: 52 x 12
##    Name    House Start_year End_year Birth_year Death_year Gender Dynasty Source
##    <chr>   <chr>      <dbl>    <dbl>      <dbl>      <dbl> <chr>  <chr>   <chr> 
##  1 Gorm d~ Gorm         936      958        908        958 M      Jellin~ https~
##  2 Toke_G~ Gorm         985      986         NA        986 M      Jellin~ https~
##  3 Harald~ Gorm        1014     1018        963         NA M      Jellin~ https~
##  4 Knud 1~ Gorm        1018     1035         NA       1035 M      Jellin~ https~
##  5 Hardek~ Gorm        1035     1042        995       1042 M      Jellin~ https~
##  6 Magnus~ Fair~       1042     1047       1018       1047 M      Jellin~ https~
##  7 Svend ~ Estr~       1047     1074       1024       1076 M      <NA>    https~
##  8 Harald~ Estr~       1074     1080       1019       1080 M      Jellin~ https~
##  9 Knud 2~ Estr~       1080     1086       1041       1086 M      Jellin~ https~
## 10 Oluf 1~ Estr~       1086     1095       1043       1095 M      Jellin~ https~
## # i 42 more rows
## # i 3 more variables: BirthDMY <chr>, DeathDMY <chr>, duration <dbl>
\end{verbatim}

\begin{Shaded}
\begin{Highlighting}[]
\NormalTok{long\_reign\_kings }\SpecialCharTok{\%\textgreater{}\%}
  \FunctionTok{count}\NormalTok{()}
\end{Highlighting}
\end{Shaded}

\begin{verbatim}
## # A tibble: 1 x 1
##       n
##   <int>
## 1    52
\end{verbatim}

We see here a list of 25 kings (and queens) who had longer than average
duration of reign.

\subsection{How many days did the three longest-ruling monarchs
rule?}\label{how-many-days-did-the-three-longest-ruling-monarchs-rule}

\begin{itemize}
\tightlist
\item
  Sort kings by reign \texttt{duration} in the descending order. Select
  the three longest-ruling monarchs with the \texttt{slice()} function
\item
  Use \texttt{mutate()} to create \texttt{Days} column where you
  calculate the total number of days they ruled
\item
  BONUS: consider the transition year (with 366 days) in your
  calculation!
\end{itemize}

\begin{Shaded}
\begin{Highlighting}[]
\NormalTok{kings\_cleaned }\SpecialCharTok{\%\textgreater{}\%}
  \FunctionTok{arrange}\NormalTok{(}\FunctionTok{desc}\NormalTok{(duration)) }\SpecialCharTok{\%\textgreater{}\%}    
  \FunctionTok{slice}\NormalTok{(}\DecValTok{1}\SpecialCharTok{:}\DecValTok{3}\NormalTok{) }\SpecialCharTok{\%\textgreater{}\%}                 
  \FunctionTok{mutate}\NormalTok{(}
    \AttributeTok{Days =}\NormalTok{ duration }\SpecialCharTok{*} \FloatTok{365.25}\NormalTok{,}
    \AttributeTok{Days =}\NormalTok{ Days}
\NormalTok{  ) }\SpecialCharTok{\%\textgreater{}\%}
  \FunctionTok{select}\NormalTok{(Name, duration, Days)   }
\end{Highlighting}
\end{Shaded}

\begin{verbatim}
## # A tibble: 3 x 3
##   Name               duration   Days
##   <chr>                 <dbl>  <dbl>
## 1 Christian 4.             60 21915 
## 2 Margrete 2.              51 18628.
## 3 Erik 7. af Pommern       43 15706.
\end{verbatim}

\subsection{Challenge: Plot the kings' duration of reign through
time}\label{challenge-plot-the-kings-duration-of-reign-through-time}

What is the long-term trend in the duration of reign among Danish
monarchs? How does it relate to the historical violence trends ?

\begin{itemize}
\tightlist
\item
  Try to plot the duration of reign column in \texttt{ggplot} with
  \texttt{geom\_point()} and \texttt{geom\_smooth()}
\item
  In order to peg the duration (which is between 1-99) somewhere to the
  x axis with individual centuries, I recommend creating a new column
  \texttt{midyear} by adding to \texttt{startYear} the product of
  \texttt{endYear} minus the \texttt{startYear} divided by two
  (\texttt{startYear\ +\ (endYear-startYear)/2}).
\item
  Now you can plot the kings dataset, plotting \texttt{midyear} along
  the x axis and \texttt{duration} along y axis
\item
  BONUS: add a title, nice axis labels to the plot and make the theme
  B\&W and font bigger to make it nice and legible!
\end{itemize}

\begin{Shaded}
\begin{Highlighting}[]
\NormalTok{kings\_cleaned }\OtherTok{\textless{}{-}}\NormalTok{ kings\_cleaned }\SpecialCharTok{\%\textgreater{}\%}
  \FunctionTok{mutate}\NormalTok{(}\AttributeTok{midyear =}\NormalTok{ Start\_year }\SpecialCharTok{+}\NormalTok{ (End\_year }\SpecialCharTok{{-}}\NormalTok{ Start\_year) }\SpecialCharTok{/} \DecValTok{2}\NormalTok{)}

\FunctionTok{ggplot}\NormalTok{(kings\_cleaned, }\FunctionTok{aes}\NormalTok{(}\AttributeTok{x =}\NormalTok{ midyear, }\AttributeTok{y =}\NormalTok{ duration)) }\SpecialCharTok{+}
  \FunctionTok{geom\_point}\NormalTok{(}\AttributeTok{color =} \StringTok{"steelblue"}\NormalTok{, }\AttributeTok{size =} \DecValTok{2}\NormalTok{, }\AttributeTok{alpha =} \FloatTok{0.7}\NormalTok{) }\SpecialCharTok{+}    
  \FunctionTok{geom\_smooth}\NormalTok{(}\AttributeTok{method =} \StringTok{"loess"}\NormalTok{, }\AttributeTok{se =} \ConstantTok{TRUE}\NormalTok{, }\AttributeTok{color =} \StringTok{"darkred"}\NormalTok{) }\SpecialCharTok{+} 
  \FunctionTok{labs}\NormalTok{(}
    \AttributeTok{title =} \StringTok{"Trend in Duration of Reign of Danish Monarchs Over Time"}\NormalTok{,}
    \AttributeTok{x =} \StringTok{"Reign Midyear"}\NormalTok{,}
    \AttributeTok{y =} \StringTok{"Duration of Reign (Years)"}
\NormalTok{  ) }\SpecialCharTok{+}
  \FunctionTok{theme\_bw}\NormalTok{(}\AttributeTok{base\_size =} \DecValTok{14}\NormalTok{)       }
\end{Highlighting}
\end{Shaded}

\begin{verbatim}
## `geom_smooth()` using formula = 'y ~ x'
\end{verbatim}

\includegraphics{DanishMonarchs_files/figure-latex/unnamed-chunk-4-1.pdf}

And to submit this rmarkdown, knit it into html. But first, clean up the
code chunks, adjust the date, rename the author and change the
\texttt{eval=FALSE} flag to \texttt{eval=TRUE} so your script actually
generates an output. Well done!

\end{document}
